\documentclass[11pt, a4paper]{article}

% Gemeinsame Style-Definitionen (TEXINPUTS muss templates/ enthalten — siehe scripts/build_tex.py)
% bfw-style.tex — gemeinsame Style-Definitionen für alle Handouts/Aufgabenhefte
% Voraussetzung: Kompilation mit xelatex (wegen fontspec)

\usepackage[ngerman]{babel}
% Babel-ngerman macht " zu einem aktiven Shorthand (z.B. für "a -> ä). In unseren
% Texten verwenden wir " als normales Zeichen — daher Shorthand komplett ausschalten.
\shorthandoff{"}
\usepackage{geometry}
\geometry{a4paper, top=2.4cm, bottom=2.4cm, left=2.3cm, right=2.3cm,
          headheight=15pt, headsep=0.7cm, footskip=1.1cm}

\usepackage{fontspec}
% Fallback-Kette: Source Sans Pro -> Calibri -> Segoe UI -> Latin Modern Sans
% (Latin Modern ist immer mit MiKTeX vorhanden)
\IfFontExistsTF{Source Sans Pro}{%
  \setmainfont{Source Sans Pro}[Ligatures=TeX]%
}{\IfFontExistsTF{Calibri}{%
  \setmainfont{Calibri}[Ligatures=TeX]%
}{\IfFontExistsTF{Segoe UI}{%
  \setmainfont{Segoe UI}[Ligatures=TeX]%
}{%
  \setmainfont{Latin Modern Sans}[Ligatures=TeX]%
}}}
\IfFontExistsTF{Source Code Pro}{%
  \setmonofont{Source Code Pro}[Scale=0.92]%
}{\IfFontExistsTF{Consolas}{%
  \setmonofont{Consolas}[Scale=0.92]%
}{%
  \setmonofont{Latin Modern Mono}[Scale=0.92]%
}}

% Gesperrte Versalien für Labels/Titelbalken (Kapitälchen-Ersatz, funktioniert
% mit jeder OpenType-Schrift — Latin Modern Sans hat keine echten Kapitälchen)
\newcommand{\bfwlabelfont}{\footnotesize\bfseries\addfontfeatures{LetterSpace=8.0}}

\usepackage{xcolor}
% Farbpalette: hell, freundlich, modern
\definecolor{bfwPrimary}{HTML}{0EA5A4}    % Petrol/Türkis - Primär
\definecolor{bfwPrimaryDark}{HTML}{0F766E} % dunkler Petrol für Headings
\definecolor{bfwAccent}{HTML}{F59E0B}     % Amber - Akzent
\definecolor{bfwSuccess}{HTML}{10B981}    % Grün - Erfolg / Lösung
\definecolor{bfwWarn}{HTML}{F97316}       % Orange - Achtung
\definecolor{bfwInk}{HTML}{0F172A}        % fast Schwarz
\definecolor{bfwInkSoft}{HTML}{475569}    % gedämpftes Grau
\definecolor{bfwBgSoft}{HTML}{F8FAFC}     % off-white
\definecolor{bfwBgTeal}{HTML}{ECFEFF}     % helles Türkis
\definecolor{bfwBgAmber}{HTML}{FEF3C7}    % helles Gelb
\definecolor{bfwBgRose}{HTML}{FFE4E6}     % helles Rosé für Eselsbrücken
\definecolor{bfwTabHead}{HTML}{E4F3F2}    % zarte Kopfzeilen-Hinterlegung für Tabellen

\usepackage[colorlinks=true, linkcolor=bfwPrimaryDark, urlcolor=bfwPrimaryDark]{hyperref}
\usepackage{lastpage}
\usepackage{enumitem}
\setlist[itemize]{leftmargin=*, itemsep=2pt, topsep=2pt}
\setlist[enumerate]{leftmargin=*, itemsep=2pt, topsep=2pt}

\usepackage[most]{tcolorbox}
\usepackage{booktabs}
\usepackage{colortbl}
\usepackage{tikz}
\usetikzlibrary{decorations.pathmorphing, positioning, shapes.geometric, arrows.meta}

% ---------------------------------------------------------------------------
% Einheitliches Tabellen-Design
%   - etwas mehr Zeilenluft, dezente graue Linien (wirkt auch mit \hline ruhiger)
%   - Kopfzeile einer Tabelle bei Bedarf zart hinterlegen: \rowcolor{bfwTabHead}
% ---------------------------------------------------------------------------
\renewcommand{\arraystretch}{1.25}
\arrayrulecolor{bfwInkSoft!50}
\setlength{\heavyrulewidth}{1pt}
\setlength{\lightrulewidth}{0.5pt}

% ---------------------------------------------------------------------------
% Boxen — klare farbige Titelbalken mit gesperrten Versal-Labels
% (Titel bleibt per Option überschreibbar: \begin{merksatz}[title={...}])
% ---------------------------------------------------------------------------
\tcbset{bfwbox/.style={%
  enhanced, breakable,
  boxrule=0.5pt, arc=2.5pt,
  left=10pt, right=10pt, top=8pt, bottom=8pt,
  toptitle=4pt, bottomtitle=4pt,
  titlerule=0pt,
  fonttitle=\bfwlabelfont,
  coltitle=white
}}

% Merkkästchen — wichtige Definition / Kernpunkt
\newtcolorbox{merksatz}[1][]{%
  bfwbox,
  colback=bfwBgTeal,
  colframe=bfwPrimaryDark,
  colbacktitle=bfwPrimaryDark,
  title={MERKE},
  #1
}

% Eselsbrücke — Mnemoniks
\newtcolorbox{eselsbruecke}[1][]{%
  bfwbox,
  colback=bfwBgRose,
  colframe=bfwAccent!85!black,
  colbacktitle=bfwAccent!85!black,
  title={ESELSBRÜCKE},
  #1
}

% Beispiel-Box
\newtcolorbox{beispiel}[1][]{%
  bfwbox,
  colback=bfwBgSoft,
  colframe=bfwInkSoft,
  colbacktitle=bfwInkSoft,
  title={BEISPIEL},
  #1
}

% Lösungs-Box (für Aufgabenhefte)
\newtcolorbox{loesung}[1][]{%
  bfwbox,
  colback=bfwSuccess!7,
  colframe=bfwSuccess!65!black,
  colbacktitle=bfwSuccess!65!black,
  title={LÖSUNG},
  #1
}

% Achtung-Box — typische Fehler / Stolperfallen
\newtcolorbox{achtung}[1][]{%
  bfwbox,
  colback=bfwWarn!6,
  colframe=bfwWarn!85!black,
  colbacktitle=bfwWarn!85!black,
  title={ACHTUNG},
  #1
}

% Formel-Box — für zentrale Formeln (groß, ruhig, ohne Titelbalken)
\newtcolorbox{formelbox}[1][]{%
  enhanced,
  colback=bfwBgTeal,
  colframe=bfwPrimaryDark,
  boxrule=1pt, arc=3pt,
  left=10pt, right=10pt, top=6pt, bottom=6pt,
  #1
}

% Glossar-Eintrag
\newcommand{\glossareintrag}[2]{%
  \par\noindent\textbf{\color{bfwPrimaryDark}#1} \enspace #2\par\smallskip
}

% ---------------------------------------------------------------------------
% Abschnitts-Design: farbiges Nummern-Quadrat + feine Linie unter \section,
% dezent farbige \subsection/\subsubsection
% ---------------------------------------------------------------------------
\usepackage{titlesec}
\newcommand{\bfwSecBadge}[1]{%
  \begin{tikzpicture}[baseline={([yshift=-0.62em]n.center)}]
    \node[fill=bfwPrimaryDark, text=white, font=\normalsize\bfseries,
          minimum width=1.55em, minimum height=1.55em, inner sep=1pt,
          rounded corners=1.5pt] (n) {#1};
  \end{tikzpicture}}
\titleformat{\section}
  {\Large\bfseries\color{bfwPrimaryDark}}
  {\bfwSecBadge{\thesection}}{0.6em}{}
  [\vspace{-0.2em}{\color{bfwPrimary!60}\titlerule[0.8pt]}]
\titleformat{\subsection}
  {\large\bfseries\color{bfwPrimaryDark}}
  {\textcolor{bfwPrimary}{\thesubsection}}{0.5em}{}
\titleformat{\subsubsection}
  {\normalsize\bfseries\color{bfwInk}}
  {\textcolor{bfwPrimary}{\thesubsubsection}}{0.4em}{}
\titlespacing*{\section}{0pt}{1.6em}{1em}
\titlespacing*{\subsection}{0pt}{1.2em}{0.5em}

% TikZ-Stil "skizzenhaft" — etwas weicher als Rough.js, aber stabil
\tikzset{
  roughbox/.style = {
    draw=bfwPrimaryDark, thick, fill=bfwBgTeal,
    rounded corners=4pt, inner sep=6pt
  },
  roughaccent/.style = {
    draw=bfwAccent, thick, fill=bfwBgAmber,
    rounded corners=4pt, inner sep=6pt
  },
  roughline/.style = {
    draw=bfwInkSoft, thick, -{Stealth[length=2.5mm]}
  }
}

% Bullet-Symbole (bewusst kein FontAwesome — vermeidet Font-Installations-Probleme)
\newcommand{\faLightbulb}{$\bullet$}
\newcommand{\faBookmark}{$\bullet$}

% ---------------------------------------------------------------------------
% Kopf-/Fußzeile
%   Kopf links:  Wortlogo „Wirtschaftsfachwirt"
%   Kopf rechts: Dokument-Kurztext, gesetzt via \kopfzeile{...}
%   Fuß links:   © Can Siebert · 2026 — Fuß rechts: Seite X von Y
%   Titelseite (Seite 1) bleibt ohne Kopf-/Fußzeile (\handouttitel setzt
%   \thispagestyle{empty}).
% ---------------------------------------------------------------------------
\usepackage{fancyhdr}
\newcommand{\bfwKopfzeileText}{}
\newcommand{\kopfzeile}[1]{\renewcommand{\bfwKopfzeileText}{#1}}
\pagestyle{fancy}
\fancyhf{}
\fancyhead[L]{\small\bfseries\color{bfwPrimaryDark}Wirtschaftsfachwirt}
\fancyhead[R]{\small\color{bfwInkSoft}\bfwKopfzeileText}
\renewcommand{\headrulewidth}{0.6pt}
\renewcommand{\headrule}{{\color{bfwPrimary}\hrule width\headwidth height 0.6pt}}
\fancyfoot[L]{\small\color{bfwInkSoft}\copyright\ Can Siebert\,\textperiodcentered\,2026}
\fancyfoot[R]{\small\color{bfwInkSoft}Seite \thepage\ von \pageref*{LastPage}}
% Auch \thispagestyle{plain} (z.B. durch \maketitle) soll leer bleiben:
\fancypagestyle{plain}{\fancyhf{}\renewcommand{\headrulewidth}{0pt}\renewcommand{\headrule}{}}

% ---------------------------------------------------------------------------
% \handouttitel{Obertitel}{Untertitel}{Kurs-Label}{Termin-Zeile}
%   Einheitlicher professioneller Titelblock: Dachzeile, farbige Headline,
%   Untertitel, farbige Trennlinie und dezente Meta-Zeile (Kurs/Termin/Dozent).
%   Ersetzt \title/\maketitle. Direkt nach \begin{document} aufrufen.
% ---------------------------------------------------------------------------
\newcommand{\handouttitel}[4]{%
  \thispagestyle{empty}%
  \begingroup
  \setlength{\parindent}{0pt}%
  {\bfwlabelfont\color{bfwPrimary}WIRTSCHAFTSFACHWIRT (IHK)\par}
  \vspace{0.55em}
  {\huge\bfseries\color{bfwPrimaryDark}#1\par}
  \vspace{0.5em}
  {\large\color{bfwInkSoft}#2\par}
  \vspace{0.9em}
  {\color{bfwPrimary}\rule{\linewidth}{2pt}}\par
  \vspace{0.55em}
  {\small\color{bfwInkSoft}#3\ \,\textperiodcentered\,\ #4\ \,\textperiodcentered\,\ Dozent: Can Siebert\par}
  \endgroup
  \vspace{1.4em}%
}


\kopfzeile{Handout VWL Lektion 2 \textperiodcentered{} Teil 1}

\begin{document}
\handouttitel{Lektion~2 \textperiodcentered{} Teil~1: Wettbewerbspolitik \& Staatseingriffe}%
  {Vom funktionierenden Wettbewerb zu Höchst- und Mindestpreisen}%
  {1.1 Volkswirtschaftliche Grundlagen}%
  {Sa, 12.09.2026 \textperiodcentered{} Session 1}

\begin{merksatz}[title={\bfseries Worum geht es in Teil 1?}]
Basis der sozialen Marktwirtschaft ist ein funktionierender Markt mit ausreichendem
Wettbewerb. Da ein sich selbst überlassener Markt die Tendenz zeigt, Wettbewerb zu
reduzieren, greift der Staat regulierend ein: durch \emph{Wettbewerbspolitik} (GWB,
Bundeskartellamt) und durch \emph{Eingriffe in die Preisbildung} (Höchst- und
Mindestpreise, Subventionen, Steuern). Beide Themen sind Standardstoff der IHK-Prüfung.
In Teil~2 (Session~2) folgt darauf die Erfolgsmessung: Volkswirtschaftliche
Gesamtrechnung und die Ziele der Stabilitätspolitik.
\end{merksatz}

\tableofcontents
\vspace{1em}
\hrule
\vspace{1em}

\section{Wettbewerb und Wettbewerbspolitik}

\subsection{Was ist Wettbewerb?}

Unter \textbf{Wettbewerb} versteht man allgemein die \emph{Rivalität der Marktteilnehmer
unter- und gegeneinander}, um ihre Ziele zulasten der Konkurrenten durchzusetzen. Der
Wettbewerb sichert u.\,a. die persönliche und wirtschaftliche Freiheit, die
Güterverteilung und Bedarfsdeckung über den Marktmechanismus und den ökonomischen
Einsatz der Produktionsfaktoren.

\subsection{Funktionen des Wettbewerbs}

Idealtypisch erfüllt der Wettbewerb folgende Funktionen:

\begin{itemize}
  \item \textbf{Freiheitsfunktion:} Jeder kann seine Fähigkeiten und Mittel für frei
    gewählte Zwecke einsetzen (Freiheit der Unternehmertätigkeit, Freiheit der
    Konsumwahl) --- freier Wettbewerb schafft mehr Handlungsalternativen.
  \item \textbf{Kontrollfunktion:} Anbieter und Nachfrager kontrollieren sich gegenseitig;
    jeder muss sich an der Leistung der Konkurrenz messen lassen. Die wirtschaftliche
    Macht des Einzelnen wird begrenzt.
  \item \textbf{Steuerungsfunktion (Lenkungsfunktion):} Mit zwei Unterfunktionen ---
    \emph{Koordinations- und Anpassungsfunktion} (Wettbewerb beschleunigt die Koordination
    von Angebot und Nachfrage; langfristig wird nur angeboten, was der Nachfrage
    entspricht) und \emph{Allokationsfunktion} (die Produktionsfaktoren Arbeit, Boden,
    Kapital, Know-how werden so kombiniert, dass ein Maximum an Produktivität entsteht).
  \item \textbf{Innovationsfunktion:} Wettbewerb schafft Anreize, Produkte zu verbessern,
    die Produktpalette zu erneuern, kostengünstiger zu produzieren und den technischen
    Fortschritt rasch zu nutzen.
  \item \textbf{Auslesefunktion:} Auslese der Marktteilnehmer nach wirtschaftlicher
    Leistungsfähigkeit --- stets kommen die leistungsfähigsten Anbieter und Nachfrager
    zum Zug.
  \item \textbf{Verteilungsfunktion:} Die Entlohnung der Produktionsfaktoren entspricht
    deren Leistung --- leistungsorientierte Einkommensverteilung.
\end{itemize}

\begin{eselsbruecke}
\textbf{F-K-S-I-A-V} --- „\textbf{F}reier \textbf{K}onkurrenzkampf \textbf{s}ichert
\textbf{i}mmer \textbf{a}lle \textbf{V}orteile``: \textbf{F}reiheit, \textbf{K}ontrolle,
\textbf{S}teuerung, \textbf{I}nnovation, \textbf{A}uslese, \textbf{V}erteilung ---
die sechs Funktionen des Wettbewerbs.
\end{eselsbruecke}

\subsection{Ziele und Instrumente der Wettbewerbspolitik}

Weil unvollkommene Märkte zu Wettbewerbsverzerrungen führen und Märkte die Tendenz haben,
Wettbewerb abzubauen, steuert der Staat im Rahmen der \textbf{Ordnungspolitik} gegen.
Schwerpunkte: freier Zugang zu Märkten, Abbau von Handelshemmnissen, Verhinderung
wettbewerbsbeschränkender Verhaltensweisen. Leitbild ist heute nicht mehr die
vollständige Konkurrenz, sondern der Erhalt von \emph{Qualität und Intensität} des
Wettbewerbs --- Wissen, Kreativität und Innovationen gelten als treibende Marktkräfte.

Die beiden gesetzlichen Grundlagen in Deutschland:

\begin{itemize}
  \item \textbf{UWG} --- Gesetz gegen den unlauteren Wettbewerb: garantiert fairen
    Wettbewerb und faire Marktpraktiken (Schutz der \emph{Art und Weise} des Wettbewerbs).
    Unzulässig sind z.\,B. irreführende und aggressive Werbung, Schleichwerbung oder die
    gezielte Herabsetzung von Mitbewerbern.
  \item \textbf{GWB} --- Gesetz gegen Wettbewerbsbeschränkungen („Kartellgesetz``):
    unterbindet Kartelle, Absprachen und Preisbindungen und wirkt der Konzentration der
    Wirtschaft entgegen (Schutz des \emph{Bestands} von Wettbewerb).
\end{itemize}

Über die Einhaltung des GWB wacht das \textbf{Bundeskartellamt}. In der EU liegt die
Wettbewerbspolitik bei der \textbf{EU-Kommission} (Kommissar für Wettbewerb);
Rechtsgrundlagen sind der \textbf{AEUV} (Vertrag über die Arbeitsweise der EU) und die
europäische \textbf{Fusionskontrollverordnung (FKVO)}. Durch die GWB-Novellen (7.~Novelle
2005: Angleichung an EU-Recht; 8.~Novelle 2013: Neustrukturierung der Missbrauchsaufsicht,
Angleichung der Fusionskontrolle; 9.~Novelle 2017: Umsetzung der
EU-Kartellschadenersatzrichtlinie) ähneln sich deutsches und europäisches Recht sehr.
Jüngere Novellen: Die 10.~Novelle (GWB-Digitalisierungsgesetz 2021) führte u.\,a.
\S\,19a GWB ein --- erweiterte Aufsicht über marktübergreifend bedeutende
Digitalkonzerne ---, die 11.~Novelle (2023) erweiterte die Befugnisse für
Sektoruntersuchungen.
\emph{EU-Recht steht über nationalem Recht} --- es greift, wenn Wettbewerbsbeschränkungen
den Handel zwischen Mitgliedstaaten beeinträchtigen. Ein weltweites Wettbewerbsrecht oder
Kartellamt existiert nicht.

\begin{merksatz}[title={\bfseries Die drei Säulen des GWB}]
\textbf{1. Kartellverbot} (\S\,1 GWB) \quad \textbf{2. Missbrauchsaufsicht}
(\S\S\,19\,ff. GWB) \quad \textbf{3. Fusionskontrolle} (\S\,35 GWB).
Merkhilfe: \textbf{Ka-Mi-Fu}.
\end{merksatz}

\begin{center}
\begin{tikzpicture}[font=\small,
    saeule/.style={rounded corners=4pt, thick, align=center, text width=3.6cm,
                   minimum height=2.9cm, inner sep=6pt}]
  \node[roughbox, minimum width=13.4cm, minimum height=1.05cm, align=center] (dach) at (0,0)
    {\bfseries GWB --- Gesetz gegen Wettbewerbsbeschränkungen („Kartellgesetz``)};
  \node[saeule, draw=bfwPrimaryDark, fill=bfwBgTeal] (s1) at (-4.7,-2.9)
    {{\color{bfwPrimaryDark}\bfseries 1. Kartellverbot}\\[1pt]
     {\footnotesize \S\,1 GWB}\\[4pt]
     {\footnotesize Preis-, Submissions-, Gebiets- und Quotenkartelle verboten;
      Ausnahmen: \S\S\,2--3 GWB, EU-Freistellungen}};
  \node[saeule, draw=bfwAccent!85!black, fill=bfwBgAmber] (s2) at (0,-2.9)
    {{\color{bfwAccent!85!black}\bfseries 2. Missbrauchsaufsicht}\\[1pt]
     {\footnotesize \S\S\,19\,ff. GWB}\\[4pt]
     {\footnotesize Ausbeutungs- und Behinderungsmissbrauch (z.\,B. Dumping)
      marktbeherrschender Unternehmen}};
  \node[saeule, draw=bfwWarn!85!black, fill=bfwBgRose] (s3) at (4.7,-2.9)
    {{\color{bfwWarn!85!black}\bfseries 3. Fusionskontrolle}\\[1pt]
     {\footnotesize \S\,35 GWB}\\[4pt]
     {\footnotesize Anmeldepflicht, Untersagung, Ministererlaubnis (\S\,42 GWB)
      als Ausnahme}};
  \node[rounded corners=4pt, thick, draw=bfwInkSoft, fill=bfwBgSoft, align=center,
        minimum width=13.4cm, minimum height=0.95cm] (amt) at (0,-5.6)
    {\bfseries Bundeskartellamt {\normalfont\footnotesize --- wacht über das GWB
     \textperiodcentered{} EU-Ebene: EU-Kommission (Art.~101 AEUV, FKVO)}};
  \draw[roughline] (s1.north) -- (s1.north |- dach.south);
  \draw[roughline] (s2.north) -- (s2.north |- dach.south);
  \draw[roughline] (s3.north) -- (s3.north |- dach.south);
  \draw[roughline] (amt.north -| s1.south) -- (s1.south);
  \draw[roughline] (amt.north -| s2.south) -- (s2.south);
  \draw[roughline] (amt.north -| s3.south) -- (s3.south);
\end{tikzpicture}
\end{center}

\subsection{Das Kartellverbot (\S\,1 GWB)}

\textbf{Kartelle} sind Kooperationen von Unternehmen auf Basis (schriftlicher oder
mündlicher) vertraglicher Vereinbarungen, um die Wettbewerbsposition zu verbessern. Die
Beteiligten bleiben \emph{rechtlich selbstständig}, schränken ihre \emph{wirtschaftliche}
Selbstständigkeit aber teilweise ein. Nach \S\,1 GWB sind Vereinbarungen zwischen
Unternehmen, Beschlüsse von Unternehmensvereinigungen und aufeinander abgestimmte
Verhaltensweisen \textbf{verboten}, wenn sie eine \emph{Verhinderung, Einschränkung oder
Verfälschung des Wettbewerbs} bewirken. Das europarechtliche Pendant ist
\textbf{Art.~101 AEUV}. Verbotene Kartellarten insbesondere:

\begin{center}
\begin{tabular}{p{4.2cm} p{9.8cm}}
\toprule
\textbf{Kartellart} & \textbf{Inhalt der Absprache}\\
\midrule
Preiskartell & Absprache einer einheitlichen Preispolitik\\
Submissionskartell & Vereinbarung von Preisuntergrenzen bzw. Festlegung, welches Kartellmitglied bei öffentlichen Ausschreibungen zum Zug kommt\\
Gebietskartell & Aufteilung des Absatzmarktes\\
Quotenkartell & künstliche Verknappung des Angebots durch Absprache von Produktionsquoten\\
\bottomrule
\end{tabular}
\end{center}

\begin{achtung}
Abgrenzung (beliebte Prüfungsfrage): Beim \textbf{Kartell} bleiben die Unternehmen
\emph{rechtlich} selbstständig und geben nur einen Teil ihrer \emph{wirtschaftlichen}
Selbstständigkeit auf. Beim \textbf{Konzern} bleiben sie rechtlich selbstständig, stehen
aber unter einheitlicher Leitung. Bei der \textbf{Fusion} verschmelzen sie auch
rechtlich zu einer Einheit --- dann greift die Fusionskontrolle (\S\,35 GWB), nicht das
Kartellverbot.
\end{achtung}

\subsubsection{Ausnahmen vom Kartellverbot}

Obwohl Kartelle grundsätzlich verboten sind, gibt es \textbf{Ausnahmen} in drei Gruppen:

\begin{enumerate}
  \item \textbf{Besondere Ausnahmeregelungen nach dem GWB:} z.\,B.
    \emph{Mittelstandskartelle} (\S\,3 GWB) --- Vereinbarungen zwischen kleinen und
    mittleren Unternehmen, die wegen ihres begrenzten Potenzials Wettbewerbsnachteile
    gegenüber Großunternehmen haben.
  \item \textbf{EU-Gruppenfreistellungsverordnungen:} Bestimmte Gruppen von Vereinbarungen
    sind freigestellt, weil sie sich trotz gewisser Wettbewerbsbeschränkung positiv auf
    den Markt auswirken; Einzelfallprüfung entfällt.
  \item \textbf{Einzelfreistellungen (Legalausnahmen, \S\,2 GWB / Art.~101 Abs.~3 AEUV):}
    Vereinbarungen sind freigestellt, wenn sie (a) den Verbraucher angemessen am Gewinn
    beteiligen, (b) der Verbesserung der Warenerzeugung oder -verteilung dienen oder
    (c) zur Förderung des technischen Fortschritts beitragen --- \emph{ohne} dass der
    Wettbewerb zwischen den Beteiligten ausgeschaltet wird.
\end{enumerate}

Nach dem \textbf{Selbstprüfungssystem} müssen die Unternehmen selbst einschätzen, ob ihr
Verhalten mit dem Kartellrecht vereinbar ist --- Anmelde- und Überprüfungsverfahren beim
Kartellamt entfallen (weniger Bürokratie, mehr Eigenverantwortung). Bei Verstößen drohen
empfindliche Strafen (Bußgelder, Gewinnabschöpfung).

\subsection{Missbrauchsaufsicht (\S\S\,19\,ff. GWB)}

Trotz Wettbewerbspolitik entstehen immer mehr Unternehmen mit erheblicher Marktmacht.
Über \emph{marktbeherrschende} Unternehmen übt das Bundeskartellamt die
\textbf{Missbrauchsaufsicht} aus. Zwei Missbrauchsformen:

\begin{itemize}
  \item \textbf{Ausbeutungsmissbrauch:} Forderung überhöhter Preise von Kunden oder
    extremer Preisnachlässe von abhängigen Zulieferern (viele Zulieferer sind so stark an
    einen Kunden gebunden, dass ihre Existenz bedroht ist, wenn sie ihn verlieren).
  \item \textbf{Behinderungsmissbrauch:} Marktmacht wird gezielt eingesetzt, um aktuelle
    oder potenzielle Konkurrenten zu behindern, zu verdrängen oder ihnen den Marktzugang
    zu versperren --- z.\,B. über \emph{Kampfpreise (Dumping)} oder die Verweigerung des
    Zugangs zu Infrastruktureinrichtungen.
\end{itemize}

\subsection{Fusionskontrolle (\S\,35 GWB)}

Während die Missbrauchsaufsicht \emph{bestehende} marktbeherrschende Unternehmen betrifft,
ist die \textbf{Fusionskontrolle} auch eine \emph{vorbeugende} Maßnahme: Geplante Fusionen
müssen beim Bundeskartellamt \emph{angemeldet} werden. Es kann den Zusammenschluss
\textbf{untersagen}, wenn die Gefahr besteht, dass eine marktbeherrschende Stellung
entsteht. Der Bundeswirtschaftsminister kann einen untersagten Zusammenschluss dennoch
erlauben (\textbf{Ministererlaubnis}, \S\,42 GWB), wenn er von gesamtwirtschaftlichem
Vorteil ist und die marktwirtschaftliche Ordnung nicht gefährdet wird.

\begin{beispiel}
Drei Bauunternehmen sprechen vor einer öffentlichen Ausschreibung ab, wer das günstigste
Angebot abgibt und den Auftrag erhält $\Rightarrow$ verbotenes
\textbf{Submissionskartell} (\S\,1 GWB). Ein marktbeherrschender Onlinehändler verkauft
gezielt unter Einstandspreis, um einen neuen Konkurrenten zu verdrängen $\Rightarrow$
\textbf{Behinderungsmissbrauch} (Dumping). Zwei große Lebensmittelketten wollen
fusionieren und hätten danach regional über 50~\% Marktanteil $\Rightarrow$ Anmeldepflicht,
das Bundeskartellamt kann die Fusion \textbf{untersagen}.
\end{beispiel}

\section{Eingriffe des Staates in die Preisbildung}

In einer Marktwirtschaft bildet sich der Preis durch Angebot und Nachfrage. In der
\emph{sozialen} Marktwirtschaft übernimmt der Staat regulierende Aufgaben, um
Benachteiligungen von Anbietern oder Nachfragern auszugleichen. Grundsätzlich gilt:

\begin{merksatz}
\textbf{Marktkonforme (indirekte) Eingriffe} beeinflussen Angebot und Nachfrage, ohne den
Marktpreismechanismus außer Kraft zu setzen (z.\,B. Steuern, Zölle, Subventionen).
\textbf{Marktkonträre (direkte) Eingriffe} greifen unmittelbar in die Preisbildung ein
und hebeln den Preismechanismus aus (z.\,B. Höchst- und Mindestpreise).
\end{merksatz}

\begin{center}
\begin{tabular}{p{3.2cm} p{5.2cm} p{5.2cm}}
\toprule
\rowcolor{bfwTabHead} & \textbf{marktkonform (indirekt)} & \textbf{marktkonträr (direkt)}\\
\midrule
Preismechanismus & bleibt in Kraft & wird außer Kraft gesetzt\\
Ansatzpunkt & Angebot bzw. Nachfrage werden verschoben & Preis wird unmittelbar staatlich festgesetzt\\
Beispiele & Steuern, Zölle, Subventionen, Transferleistungen, Mengenmaßnahmen & Höchstpreis, Mindestpreis\\
Typische Folgen & Preissignale bleiben erhalten (aber Verzerrung bei Dauersubventionen) & dauerhafte Überschüsse bzw. Knappheiten, graue und schwarze Märkte, Folgekosten\\
\bottomrule
\end{tabular}
\end{center}

\subsection{Indirekte (marktkonforme) Maßnahmen}

\begin{itemize}
  \item \textbf{Verbrauchssteuern:} Höhere Verbrauchssteuern (z.\,B. Tabaksteuer)
    verteuern das Produkt $\Rightarrow$ höherer Gleichgewichtspreis, geringere
    Gleichgewichtsmenge. So kann Verbraucherverhalten aus umwelt- und
    gesundheitspolitischen Gründen gelenkt werden (weniger Raucher). Je nach
    Nachfragereaktion können die Steuereinnahmen trotz höheren Steuersatzes sinken ---
    bei unelastischer Nachfrage steigen sie sogar. Steuersenkungen wirken umgekehrt:
    niedrigerer Gleichgewichtspreis, höhere Menge.
  \item \textbf{Zölle:} Senkung von Einfuhrzöllen erhöht die Importe (wirkt gegen
    Preissteigerungen im Inland); Erhöhung von Einfuhrzöllen verteuert Importwaren und
    schützt die einheimische Industrie. Die Gesetzgebungs- und Ertragskompetenz für Zölle
    liegt bei der \emph{EU}.
  \item \textbf{Subventionen und Steuervergünstigungen (Anbieterförderung):} Staatliche
    Finanzhilfen ermöglichen es Anbietern, zu einem Preis unterhalb der
    Produktionskosten zu verkaufen --- z.\,B. zur Markteinführung neuer Produkte
    (Windkraftanlagen). \emph{Problematisch bei Dauersubvention:} Sie behindert sämtliche
    Preisfunktionen; der künstlich niedrige Preis sendet falsche Signale, das Ausscheiden
    kostenungünstiger Anbieter wird verhindert, Produktionsfaktoren bleiben in nicht
    (mehr) gefragten Produkten gebunden.
  \item \textbf{Transferleistungen (Subvention der Nachfrager):} Der Staat gibt Haushalten
    zusätzliche Mittel, damit sie ein bestimmtes Gut erwerben können --- z.\,B.
    \emph{Wohngeld} oder \emph{Kindergeld}. Die Finanzierung von Subventionen und
    Transfers tragen die Steuerzahler insgesamt.
  \item \textbf{Mengenmaßnahmen:} Regulierung der Angebots- und Nachfragemenge, z.\,B.
    Einfuhrkontingente, Einfuhrverbote, staatliche Vorratshaltung.
\end{itemize}

\subsection{Direkte (marktkonträre) Maßnahmen: Mindest- und Höchstpreise}

Beide Preisgrenzen lassen sich im Preis-Mengen-Diagramm darstellen --- der staatlich
festgesetzte Preis „klemmt`` den Markt außerhalb des Gleichgewichts fest:

\begin{center}
\begin{tikzpicture}[font=\footnotesize, scale=0.94]
  % ---------- linkes Diagramm: Mindestpreis ----------
  \begin{scope}
    \node[font=\small\bfseries, color=bfwPrimaryDark, align=center] at (2.7,5.3)
      {Mindestpreis\\[-1pt]{\footnotesize\color{bfwInkSoft}\mdseries Schutz der Anbieter}};
    \draw[-{Stealth[length=2.5mm]}, semithick, color=bfwInk] (0,0) -- (0,4.55) node[above] {$p$};
    \draw[-{Stealth[length=2.5mm]}, semithick, color=bfwInk] (0,0) -- (5.3,0) node[right] {$x$};
    \draw[bfwPrimaryDark, very thick] (0.4,0.5) -- (4.4,3.9) node[above] {$A$};
    \draw[bfwWarn!90!black, very thick] (0.4,3.9) -- (4.4,0.5) node[right] {$N$};
    \draw[densely dashed, color=bfwInkSoft] (0,2.2) -- (2.4,2.2) -- (2.4,0);
    \node[left, color=bfwInkSoft] at (0,2.2) {$p_0$};
    \node[below, color=bfwInkSoft] at (2.4,0) {$x_0$};
    \fill[bfwInk] (2.4,2.2) circle (1.7pt);
    \draw[bfwAccent, line width=1.5pt] (0,3.05) -- (5.05,3.05);
    \node[left, color=bfwAccent!80!black] at (0,3.05) {$p_{\mathrm{min}}$};
    \draw[densely dashed, color=bfwInkSoft] (1.4,3.05) -- (1.4,0);
    \draw[densely dashed, color=bfwInkSoft] (3.4,3.05) -- (3.4,0);
    \node[below, color=bfwInkSoft] at (1.4,0) {$x_N$};
    \node[below, color=bfwInkSoft] at (3.4,0) {$x_A$};
    \fill[bfwAccent!80!black] (1.4,3.05) circle (1.5pt) (3.4,3.05) circle (1.5pt);
    \draw[{Stealth[length=2mm]}-{Stealth[length=2mm]}, color=bfwAccent!80!black, semithick]
      (1.4,3.5) -- (3.4,3.5)
      node[midway, above, align=center] {\bfseries Angebots-\\[-2pt]\bfseries überschuss};
  \end{scope}
  % ---------- rechtes Diagramm: Höchstpreis ----------
  \begin{scope}[xshift=7.5cm]
    \node[font=\small\bfseries, color=bfwWarn!85!black, align=center] at (2.7,5.3)
      {Höchstpreis\\[-1pt]{\footnotesize\color{bfwInkSoft}\mdseries Schutz der Nachfrager}};
    \draw[-{Stealth[length=2.5mm]}, semithick, color=bfwInk] (0,0) -- (0,4.55) node[above] {$p$};
    \draw[-{Stealth[length=2.5mm]}, semithick, color=bfwInk] (0,0) -- (5.3,0) node[right] {$x$};
    \draw[bfwPrimaryDark, very thick] (0.4,0.5) -- (4.4,3.9) node[above] {$A$};
    \draw[bfwWarn!90!black, very thick] (0.4,3.9) -- (4.4,0.5) node[right] {$N$};
    \draw[densely dashed, color=bfwInkSoft] (0,2.2) -- (2.4,2.2) -- (2.4,0);
    \node[left, color=bfwInkSoft] at (0,2.2) {$p_0$};
    \node[below, color=bfwInkSoft] at (2.4,0) {$x_0$};
    \fill[bfwInk] (2.4,2.2) circle (1.7pt);
    \draw[bfwAccent, line width=1.5pt] (0,1.35) -- (5.05,1.35);
    \node[left, color=bfwAccent!80!black] at (0,1.35) {$p_{\mathrm{max}}$};
    \draw[densely dashed, color=bfwInkSoft] (1.4,1.35) -- (1.4,0);
    \draw[densely dashed, color=bfwInkSoft] (3.4,1.35) -- (3.4,0);
    \node[below, color=bfwInkSoft] at (1.4,0) {$x_A$};
    \node[below, color=bfwInkSoft] at (3.4,0) {$x_N$};
    \fill[bfwAccent!80!black] (1.4,1.35) circle (1.5pt) (3.4,1.35) circle (1.5pt);
    \draw[{Stealth[length=2mm]}-{Stealth[length=2mm]}, color=bfwAccent!80!black, semithick]
      (1.4,0.75) -- (3.4,0.75)
      node[midway, fill=white, inner sep=1.5pt, align=center]
      {\bfseries Nachfrage-\\[-2pt]\bfseries überschuss};
  \end{scope}
\end{tikzpicture}
\end{center}

\begin{beispiel}
\textbf{Durchgerechnet:} Auf einem Markt gelte die Nachfragefunktion $x_N = 100 - 2p$
und die Angebotsfunktion $x_A = -20 + 2p$ ($p$ in €, $x$ in Stück).\\[0.3em]
\textbf{Schritt 1 --- Gleichgewicht:} $100 - 2p = -20 + 2p \Rightarrow 4p = 120
\Rightarrow p_0 = 30$~€; eingesetzt: $x_0 = 100 - 60 = 40$ Stück.\\
\textbf{Schritt 2 --- Mindestpreis 35~€ (über $p_0$):} $x_N = 100 - 70 = 30$,
$x_A = -20 + 70 = 50$ $\Rightarrow$ \textbf{Angebotsüberschuss von 20 Stück}
($x_N < x_0 < x_A$).\\
\textbf{Schritt 3 --- Höchstpreis 25~€ (unter $p_0$):} $x_N = 100 - 50 = 50$,
$x_A = -20 + 50 = 30$ $\Rightarrow$ \textbf{Nachfrageüberschuss von 20 Stück}
($x_A < x_0 < x_N$).
\end{beispiel}

\subsubsection{Mindestpreis}

Ein \textbf{Mindestpreis} ist eine staatlich verordnete \emph{Preisuntergrenze}, die
nicht unterschritten werden darf; er liegt \emph{über} dem Gleichgewichtspreis.
\textbf{Ziel: Schutz der Anbieter.} Folgen: Das Angebot wird ausgeweitet, die Nachfrage
sinkt $\Rightarrow$ \textbf{Angebotsüberschuss} ($x_N < x_0 < x_A$). Die Überschüsse
stellen den Staat vor neue Probleme --- er muss mengenregulierend eingreifen
(Kontingente, Prämien für Angebotsreduzierung, Flächenstilllegung). Bei
Mengenkontingenten kann ein \textbf{grauer Markt} entstehen, auf dem die Mehrproduktion
unter dem festgesetzten Mindestpreis angeboten wird.

\begin{beispiel}
\textbf{EU-Agrarmarkt:} Landwirten wurden jahrzehntelang Mindestpreise mit
Abnahmegarantien zugestanden --- Folge: massive Überproduktion („Butterberge``,
„Milchseen``). Der Staat musste die Überschüsse durch Einlagerung, Verarbeitung,
Veräußerung oder Vernichtung beseitigen (erhebliche Kosten) und das Angebot durch
Produktionsquoten und Stilllegungsprämien eindämmen.
\textbf{Gesetzlicher Mindestlohn:} Er hebelt die Findung des Gleichgewichtslohns aus und
kann zu einem dauerhaften Angebotsüberhang an Arbeitskräften (Arbeitslosigkeit) führen,
wenn Arbeitgeber Arbeitskräfte durch Technik substituieren oder Arbeitsumfänge ins
Ausland verlagern.
\end{beispiel}

\subsubsection{Höchstpreis}

Ein \textbf{Höchstpreis} ist eine staatlich verordnete \emph{Preisobergrenze}; er liegt
\emph{unter} dem Gleichgewichtspreis. \textbf{Ziel: Schutz der Nachfrager.} Beispiele:
Mietpreisbindung im sozialen Wohnungsbau, \emph{Mietpreisbremse} auf angespannten
Wohnungsmärkten. Folgen: Die Nachfrage steigt, das Angebot sinkt $\Rightarrow$
\textbf{Nachfrageüberschuss} ($x_A < x_0 < x_N$). Damit der Höchstpreis wirkt, sind
weitere Maßnahmen nötig: \emph{Subventionen}, damit überhaupt ein Angebot vorhanden ist,
sowie \emph{Kontingentierung und Rationierung} der knappen Menge. Häufig bilden sich
\textbf{Schwarzmärkte}, auf denen die Güter zu höheren Preisen verkauft werden. Selten
kann das verminderte Angebot über \emph{Importe} gelöst werden, wenn ausländische
Anbieter aufgrund von Kostenvorteilen unter dem Höchstpreis anbieten können.

\begin{eselsbruecke}
\textbf{Mi}ndestpreis liegt \textbf{über}, \textbf{Hö}chstpreis liegt \textbf{unter} dem
Gleichgewichtspreis --- sonst wären sie wirkungslos! Mindestpreis schützt die
\textbf{A}nbieter $\Rightarrow$ \textbf{A}ngebotsüberschuss; Höchstpreis schützt die
\textbf{N}achfrager $\Rightarrow$ \textbf{N}achfrageüberschuss. Merkhilfe: \emph{Die
geschützte Marktseite weitet ihre Menge aus --- auf ihrer Seite entsteht der Überschuss.}
\end{eselsbruecke}

\begin{merksatz}[title={\bfseries Beurteilung staatlicher Preiseingriffe (prüfungsrelevant!)}]
Staatliche Preisfestsetzungen setzen die Funktionen des Marktpreises (Signal, Ausgleich,
Allokation, Selektion) außer Kraft: Es entstehen dauerhafte Überschüsse bzw. Knappheiten,
Folgekosten für den Staat, Fehlallokation von Produktionsfaktoren sowie graue und schwarze
Märkte. Marktkonforme Eingriffe (Steuern, Subventionen, Transfers) sind daher
ordnungspolitisch vorzuziehen --- aber auch Dauersubventionen verzerren die Preissignale.
\end{merksatz}

\begin{merksatz}[title={\bfseries Wichtig für die Prüfungsvorbereitung (Teil 1)}]
Sie sollten sicher können: die \emph{Funktionen des Wettbewerbs} erläutern, die
\emph{drei Säulen des GWB} mit Kartellarten, Ausnahmen und Beispielen erklären,
\emph{marktkonforme und marktkonträre} Eingriffe unterscheiden sowie Höchst- und
Mindestpreise \emph{im Preis-Mengen-Diagramm} einzeichnen und ihre Folgen
(Überschüsse, graue und schwarze Märkte, Folgekosten) beurteilen. Üben Sie mit dem
Aufgabenheft (Aufgabe~4: Mindestpreis) und dem Quiz zu Teil~1!
\end{merksatz}

\vspace{1em}
\hrule
\vspace{0.8em}
\noindent\textsf{\small Quellen: Rahmenplan Wirtschaftsbezogene Qualifikationen (DIHK), Abschnitte 1.1.1.2, 1.1.1.3;
Textband Volkswirtschaft (DIHK-Bildungs-GmbH), Kap.~1.2--1.3 (S.~12--18).
Weiter geht es in Session~2 (Teil~2): Volkswirtschaftliche Gesamtrechnung \& Ziele der Stabilitätspolitik.}

\end{document}
